\begin{exercise}{3}
Show that every map in $\Delta$ can be uniquely factored as a
composition of $s_i$'s followed by a composition of $d_j$'s. Therefore,
a simplicial set is equivalently described by a sequence of sets $X_n$
equipped with face and degeneracy maps satisfying the simplicial
identities.

\begin{proof}
For $\alpha\colon[n]\to[m]$, write its image as
$\{c_0<\cdots<c_r\}$. There is a unique factorization
\[
 [n]\xrightarrow{\rho}[r]\xrightarrow{\iota}[m],\qquad
 \rho(t)=k\ \text{if }\alpha(t)=c_k,\quad \iota(k)=c_k,
\]
with $\rho$ surjective and $\iota$ injective.
Let $a_1<\cdots<a_p$ list the values omitted by $\alpha$, and let
$b_1<\cdots<b_q$ list the positions $b<n$ with
$\alpha(b)=\alpha(b+1)$. Inserting the omitted values in increasing
order and merging the repeated positions in decreasing order gives
\begin{equation}\label{eq:delta-normal-form}
 \alpha=d_{a_p}\cdots d_{a_1}s_{b_1}\cdots s_{b_q},
 \qquad p=m-r,\quad q=n-r.
\end{equation}
Empty products mean identity maps. The two index lists are determined
by $\alpha$, proving uniqueness in this ordered normal form.
An ordering convention is essential: arbitrary words are not unique,
as $s_0s_0=s_0s_1$ shows.

We also need that the identities in Exercise~\ref{ex:2} account for
all relations between these generators. Read those identities in
$\Delta$. The mixed identities move each $s_jd_i$ pair into $ds$ order
or cancel it; each step decreases the number of pairs consisting of
an $s$ to the left of a $d$. Within the resulting pure blocks, use
\[
 d_id_j\longmapsto d_{j+1}d_i\quad(i\le j),\qquad
 s_js_i\longmapsto s_is_{j+1}\quad(i\le j).
\]
The face index sequence strictly increases lexicographically, and the
degeneracy index sequence strictly increases when read from right to
left. Since the possible index sequences at fixed levels form finite
sets, these steps terminate in the ordered form
\eqref{eq:delta-normal-form}. Thus any two words for the same map are
related by the identities of Exercise~\ref{ex:2}.

Consequently, given sets $X_n$ and face and degeneracy maps satisfying
those identities, define $X(\alpha)$ by applying the corresponding
operators in reverse order in \eqref{eq:delta-normal-form}.
The reduction just proved preserves this operator composite, so
concatenation of words gives
$X(\beta\alpha)=X(\alpha)X(\beta)$; the empty word gives the identity.
This defines a simplicial set. Conversely, Definition~1.1.7 and
Exercise~\ref{ex:2} supply such operators from any simplicial set.
The constructions are inverse; morphisms are precisely the degreewise
maps commuting with all face and degeneracy operators.
\end{proof}
\end{exercise}
